\documentclass[10pt,a4paper]{article}
\usepackage[spanish,es-nodecimaldot]{babel}
\usepackage[utf8]{inputenc}
\usepackage[T1]{fontenc}
\usepackage{lmodern}
\usepackage[top=1.1cm,bottom=1.1cm,left=1.5cm,right=1.5cm]{geometry}
\usepackage{amsmath,amssymb}
\usepackage[table]{xcolor}
\usepackage{enumitem}
\usepackage[normalem]{ulem}

\setlength{\parindent}{0pt}
\setlength{\parskip}{2pt}
\pagestyle{empty}

\AtBeginDocument{%
  \setlength{\abovedisplayskip}{3pt}%
  \setlength{\belowdisplayskip}{3pt}%
  \setlength{\abovedisplayshortskip}{2pt}%
  \setlength{\belowdisplayshortskip}{2pt}%
}

\begin{document}

{\Large\bfseries Prueba U} {\large (2 muestras indep.)}
\vspace{4pt}

Existen dos pruebas: la U y la H.

\textbf{Prueba U (Wilcoxon) - Alternativa no param. de t bimuestral}

Se quiere probar si las medias de dos muestras son significativamente iguales. Se busca probar si las muestras vienen de poblaciones idénticas (Misma distribución: misma forma, $\mu$ y varianza)

\textbf{Pasos} \hspace{1.2cm} qnorm para las absicas

1.Unimos las muestras y las ordenamos de menor a mayor. (Rango=Sum de la posic.)

2. A cada valor se le indica a que muestra pertenece (I, II) y se le asigna un rango (posición). Si dos valores son iguales, se asigna un promedio de sus rangos.

3. Sumar los rangos de cada muestra R1 y R2. se comparan y se elije el Ri menor.

4. Se define el estadístico:
\[
U_i = n_1.\, n_2 + \frac{n_i.(n_i+1)}{2} \; - \; R_i
\]

\begin{minipage}{0.52\textwidth}
Si las $n_1$ y $n_2$ son $> 8$, $U_i$ tiene una dist. normal. \\
gralm. prueba 2 colas
\end{minipage}
\begin{minipage}{0.45\textwidth}
\[
z = \frac{U_i - u_{U_i}}{\sigma_{U_i}}
\]
\end{minipage}

\[
\mu_{U1} = \frac{n_1\cdot n_2}{2}
\qquad\qquad
\sigma_{U1} = \sqrt{\frac{n_1\cdot n_2\cdot(n_1+n_2+1)}{12}}
\]

\end{document}
